\documentclass[10pt,a4paper]{amsart}
\usepackage[margin=19mm]{geometry}
\usepackage{amsmath,amssymb,mathtools}
\usepackage[scr=rsfs,cal=euler]{mathalfa}
\usepackage{xfrac}
\usepackage[unicode,hidelinks,bookmarks=false]{hyperref}
\newcommand{\ie}{i.e.\ }
\newcommand{\cf}{cf.\ }
\newcommand{\resp}{respectively\ }
\newcommand{\Subsection}{\subsection}
\providecommand{\nobreakdash}{\nobreak\hbox{-}}
\setlength{\emergencystretch}{2em}
\multlinegap0pt
\usepackage{bm}
\usepackage{bbm}
\usepackage{dsfont}
\usepackage{xspace}
\usepackage{cancel}
\usepackage{mathtools}
\usepackage{amsthm}
\makeatletter
\@ifundefined{thm}{\newtheorem{thm}{Thm}}{}
\@ifundefined{prop}{\newtheorem{prop}[thm]{Proposition}}{}
\makeatother
\newcommand{\FF}{\mathbb{F}}
\newcommand{\calR}{\mathcal{R}}
\newcommand{\rmH}{\mathrm H}%for cohomology groups
\title[A cohomological translation of the Kaplansky radical for pro]{A cohomological translation of the Kaplansky radical for profinite groups}
\author{Simone Blumer, Julian Feuerpfeil, Lucas Correa Lopes, Claudio Quadrelli}
\date{Source: arXiv:2606.31547v2 (CC BY 4.0)}
\begin{document}
\maketitle

\noindent\emph{Source and licence.} A cohomological translation of the Kaplansky radical for profinite groups. Version arXiv:2606.31547v2 (CC BY 4.0), distributed under the Creative Commons Attribution 4.0 International licence, as recorded in the arXiv licence metadata for this exact version and shown on the abstract page https://arxiv.org/abs/2606.31547v2. Full rights notice: licenses/arxiv-2606.31547-CC-BY-4.0.txt.\par\smallskip

\noindent\textit{Editorial scope.} Counted unit: the Prop whose source label is \texttt{prop:D-RAAGs have KN} in the article (source0.tex lines 1274--1277, source label \texttt{prop:D-RAAGs have KN}), delivered together with the article's own complete original proof of it (source0.tex lines 1278--1299, one \texttt{proof} environment of 22 lines). No mathematics is written, repaired or re-derived: statement and proof are the article's own text line for line. Required context retained uncounted: prop:almost direct product (lines 939--942). Every definition, display and label of those lines is retained unchanged. Omitted from this package and not redistributed: 1--938, 943--1273, 1300--1385. Nothing inside a retained excerpt is omitted or altered: every retained body is delivered exactly as the article's lines are written, so no omission receipt of any kind accompanies this package. Citations: the retained excerpts carry no citation command at all, so no bibliography entry of the article is transcribed and no reference is redistributed. Reference adaptation: none was needed and none was made; every reference inside the delivered unit resolves inside it, so no unresolved reference or citation is produced. Printed slips kept literal: no printed word, symbol, hypothesis or number of the counted statement or of its proof was altered. Layout and class: the article's own class is \texttt{amsart} with packages including amsaddr, subfig, graphicx, amssymb, amsfonts, amsrefs, hyperref, tikz-cd, enumitem, array, changes, ulem, xy, bbm; the wrapper is the lane's pinned \texttt{amsart} editorial wrapper at 10pt on A4, which restates the theorem-like environments the retained text needs and adds the article's own macro definitions that the retained text needs (\texttt{\textbackslash FF}, \texttt{\textbackslash calR}, \texttt{\textbackslash rmH}), with extra wrapper packages (bm, bbm, dsfont, xspace, cancel, mathtools). The delivered numbering is the unit's own. Assets: no retained span contains a figure environment, a graphics-inclusion command, a PDF-page inclusion, a table or a TikZ picture; no image file is copied into this package, the package declares no asset dependency and no image file is redistributed with it. Licence: version arXiv:2606.31547v2 of this article is distributed under the Creative Commons Attribution 4.0 International licence, as recorded in the arXiv licence metadata for this exact version and shown on the abstract page https://arxiv.org/abs/2606.31547v2. A full rights notice is delivered at licenses/arxiv-2606.31547-CC-BY-4.0.txt. Characters: every retained excerpt is pure ASCII, so the delivered engine, XeLaTeX with its default Latin Modern text font, reports no missing character.
\begin{prop}
\label{prop:almost direct product}
    Let $G=N\rtimes Q$ be a $p$-almost direct product of profinite groups such that $\rmH^1(Q,\FF_p)$ and $\rmH^1(N,\FF_p)$ are non-trivial, then $\calR_p(G)=0$.
\end{prop}

\begin{prop}
\label{prop:D-RAAGs have KN}
    If $\Gamma$ is a non-empty finite graph, and $G_{\Gamma,\bf z}$ is a $\Delta$-RAAG defined on $\Gamma$, then $\calR_2(G_{\Gamma,\bf z})=0$, and hence $G_{\Gamma,\bf z}$ satisfies the $2$-Kijima--Nishi property.
\end{prop}

\begin{proof}
{%\color{julian} 
Since $\Delta$ acts by conjugation with the elements $\mathbf{z}_v$ on the generators $v$ of $G_\Gamma$ we see that the induced action on $\rmH^1(G_\Gamma,\FF_2)$ is trivial and hence $G_{\Gamma,\mathbf{z}}$ splits as a $2$-almost direct product and, by Proposition~\ref{prop:almost direct product}, its $\FF_2$-cup radical is trivial.}
% {\color{gray} 
% \comment{as $\Delta$-RAAGs split als $2$-almost direct product, we get this without explicit computation.}
% For all $v\in V$, there exist uniquely determined scalars $\alpha_{v,w}\in \F_2$ such that $\bar{\bf z}_v=\sum_{w\in V}\alpha_{v,w}w$, and $\alpha_{v,w}=0$, if either $\{v,w\}\in E$, or $v=w$. Here $\bar{\bf z}_v$ denotes the image of ${\bf z}_v$ into the Frattini quotient $F_{\Gamma}/\Phi(F_{\Gamma})$, where $F_\Gamma$ is the free pro-$p$ group generated by the vertices of $\Gamma$. 
% The low-degree $\F_2$-cohomology groups of $G_{\Gamma,\bf z}$ can be described as follows:
% \begin{enumerate}
%     \item $\rmH^1(G_{\Gamma,\bf z},\F_2)$ has a basis $(x_0^\ast,v^\ast)_{v\in V}$, where $v^\ast:V\to \F_2$ is the dual basis element to the canonical generator $v\in V$, and $x_0^\ast(v)=0$, $x_0^\ast(x_0)=1$;
%     \item $\rmH^2(G_{\Gamma,\bf z},\F_2)$ has a basis $(\xi,\varepsilon_v,\varepsilon_e)_{v\in V,e\in E}$.
% \end{enumerate}
% With respect to these bases, the cup product can be described by \cite[Prop.~9.3.13]{NSW2008}.
% \begin{enumerate}
%     \item $x_0^\ast\smallsmile x_0^\ast=\xi$,
%     \item $v^\ast\smallsmile v^\ast=\varepsilon_v$,
%     \item $x_0^\ast\smallsmile v^\ast=\varepsilon_v$,
%     \item $v^\ast\smallsmile w^\ast=\alpha_{v,w}\varepsilon_v$, for $\alpha_{v,w}\neq 0$,
%     \item $v^\ast\smallsmile w^\ast=\varepsilon_e$, if $\{v,w\}=e\in E$.
% \end{enumerate}
% It follows that if $0\neq\alpha\in \rmH^1(G_{\Gamma,\bf z},\F_2)$, then $\alpha\smallsmile x_0^\ast\neq 0\in \rmH^2(G_{\Gamma,\bf z},\F_2)$, proving that $\calR_p(G_{\Gamma,\bf z})=0$.
% }
\end{proof}

\end{document}
